%% Recorte vectorial de una familia de metodos, generado a partir de
%% manuscript/presentation/figures/fig_model_ladder_tikz.tex (NO modificado).
%% Fuente compartida (colores, geometria, macros) copiada tal cual; solo se aisla
%% un panel y se traducen las tres etiquetas de fila.
\definecolor{vgTarget}{HTML}{2A7F8E}  % observed target (teal)
\definecolor{vgRecon}{HTML}{1B9E77}   % reconstruction / smoother output (green)
\definecolor{vgExt}{HTML}{3B6FA0}     % external predictors / features (blue)
\definecolor{vgEdge}{HTML}{D98A2B}    % gap-edge / residual correction (amber)
\definecolor{vgPrior}{HTML}{7E6CA8}   % TS-ICL / pretrained sequence model (violet)
\definecolor{vgHidden}{HTML}{C0504D}  % hidden gap (red)
\definecolor{vgInk}{HTML}{2B2B2B}     % text
\definecolor{vgFaint}{HTML}{9AA6B2}   % row tags / faint rules
\definecolor{vgGrid}{HTML}{D9DEE2}    % daily time grid
\definecolor{vgArrow}{HTML}{8A9299}   % flow arrows
\definecolor{vgPanel}{HTML}{F7F8F9}   % panel background
\definecolor{vgChip}{HTML}{ECEFF2}    % model chip fill

\definecolor{mBase}{HTML}{5F7F8A}     % simple baselines
\definecolor{mGP}{HTML}{1B9E77}       % Gaussian process
\definecolor{mExt}{HTML}{3B6FA0}      % external tabular
\definecolor{mGap}{HTML}{D98A2B}      % gap-edge residual
\definecolor{mTS}{HTML}{7E6CA8}       % TS-ICL
\definecolor{tsgreen}{HTML}{08A69B}
\definecolor{tsblue}{HTML}{004CF8}


\begin{tikzpicture}[x=0.38cm, y=0.38cm, font=\small,
    flow/.style={-{Latex[length=2mm]}, vgArrow, line width=0.9pt}]

  % ---------------------------------------------------------------------------
  % Global geometry
  % ---------------------------------------------------------------------------
  \def\PW{9.7}        % panel width
  \def\PH{24.9}       % panel height
  \def\DX{10.55}      % horizontal shift between panels

  \def\Ytitle{23.95}
  \def\YmodelA{20.60}
  \def\YmodelB{22.70}
  \def\YhowA{10.45}
  \def\YhowB{19.75}
  \def\YoutA{1.55}
  \def\YoutB{9.20}

  \def\Yeq{7.65}      % equation height inside output box
  \def\ob{2.95}       % baseline of output mini-plots

  % ---------------------------------------------------------------------------
  % Helpers
  % ---------------------------------------------------------------------------

  \newcommand{\panelbg}[2]{%
    \fill[vgPanel, rounded corners=3pt] (-0.25,0.75) rectangle (\PW,\PH);
    \draw[#2, opacity=0.45, line width=0.65pt, rounded corners=3pt]
      (-0.25,0.75) rectangle (\PW,\PH);
    \node[font=\small\bfseries, text=vgInk, align=center] at (4.72,\Ytitle) {#1};
  }

  \newcommand{\modelchip}[1]{%
    \fill[vgChip, rounded corners=3pt] (1.05,\YmodelA) rectangle (8.40,\YmodelB);
    \draw[#1, opacity=0.40, line width=0.55pt, rounded corners=3pt]
      (1.05,\YmodelA) rectangle (8.40,\YmodelB);
  }

  \newcommand{\howbox}[1]{%
    \fill[#1, opacity=0.040, rounded corners=3pt] (0.55,\YhowA) rectangle (8.95,\YhowB);
    \draw[#1, opacity=0.75, line width=0.7pt, rounded corners=3pt]
      (0.55,\YhowA) rectangle (8.95,\YhowB);
  }

  \newcommand{\outbox}[1]{%
    \fill[#1, opacity=0.040, rounded corners=3pt] (0.55,\YoutA) rectangle (8.95,\YoutB);
    \draw[#1, opacity=0.75, line width=0.7pt, rounded corners=3pt]
      (0.55,\YoutA) rectangle (8.95,\YoutB);
  }

  \newcommand{\imageorfallback}[3]{%
    \IfFileExists{#1}{%
      \includegraphics[width=#2]{#1}%
    }{%
      \fbox{\parbox[c][2.8cm][c]{#2}{\centering\scriptsize #3}}%
    }%
  }

  % output grid with observed context and hidden interval
  \newcommand{\miniout}[1]{%
    \foreach \s in {1,...,9}{
      \draw[vgGrid,line width=0.28pt](0.2+\s*0.9,#1)--(0.2+\s*0.9,#1+3.3);
    }
    \fill[vgHidden,opacity=0.12](0.2+3.5*0.9,#1)rectangle(0.2+6.5*0.9,#1+3.3);
    \draw[vgTarget,line width=1.0pt](1.1,#1+1.35)--(2.0,#1+1.00)--(2.9,#1+1.20);
    \draw[vgTarget,line width=1.0pt](6.5,#1+2.50)--(7.4,#1+2.85)--(8.3,#1+2.50);
    \foreach \s/\y in {1/1.35,2/1.00,3/1.20,7/2.50,8/2.85,9/2.50}{
      \fill[vgTarget](0.2+\s*0.9,#1+\y)circle(0.10);
    }
  }

  % compact target-gap sketch used in "how" boxes
  \newcommand{\howgap}[1]{%
    \foreach \s in {1,...,9}{
      \draw[vgGrid,line width=0.25pt](0.2+\s*0.9,#1)--(0.2+\s*0.9,#1+3.3);
    }
    \fill[vgHidden,opacity=0.12](0.2+3.5*0.9,#1)rectangle(0.2+6.5*0.9,#1+3.3);
    \draw[vgTarget,line width=0.95pt](1.1,#1+1.35)--(2.0,#1+1.00)--(2.9,#1+1.20);
    \draw[vgTarget,line width=0.95pt](6.5,#1+2.50)--(7.4,#1+2.85)--(8.3,#1+2.50);
    \foreach \s/\y in {1/1.35,2/1.00,3/1.20,7/2.50,8/2.85,9/2.50}{
      \fill[vgTarget](0.2+\s*0.9,#1+\y)circle(0.09);
    }
  }


  \node[rotate=90,font=\footnotesize\bfseries,text=vgFaint] at (-2.25,22.10) {MÉTODO};
  \node[rotate=90,font=\footnotesize\bfseries,text=vgFaint] at (-2.25,15.10) {MECANISMO};
  \node[rotate=90,font=\footnotesize\bfseries,text=vgFaint] at (-2.25,5.35) {SALIDA};

  % ---------------------------------------------------------------------------
  \begin{scope}[shift={(0,0)}]
    \panelbg{Líneas de base simples}{mBase}
    \modelchip{mBase}
    \node[font=\tiny, text=vgInk, align=center, text width=3.2cm] at (4.72,21.65)
      {climatología\\persistencia\\interpolación lineal};

    \howbox{mBase}
    % three rule sketches with equations, not names
    \def\yshift{-0.9}
\foreach \yy/\eqn in {
  17.90/{{$\hat y_t = y_{t^-}$}},
  15.90/{{$\hat y_t=\tilde y_t^{\mathrm{lin}}$}},
  13.90/{{$\hat y_t = \mu_{\mathrm{season}}$}}
}{
  \fill[vgHidden,opacity=0.12](1.35,\yy-0.55+\yshift) rectangle (2.30,\yy+0.55+\yshift);
  \fill[vgTarget](0.95,\yy-0.25+\yshift) circle(0.09);
  \fill[vgTarget](2.80,\yy+0.30+\yshift) circle(0.09);
  \node[font=\scriptsize,text=vgInk,anchor=west] at (3.60,\yy+\yshift) {\eqn};
}
    % persistence
    \draw[mBase,line width=1.35pt] (0.95,16.75)--(2.80,16.75);
    % interpolation
    \draw[mBase,line width=1.35pt] (0.95,14.75)--(2.80,15.30);
    % climatology
    \draw[mBase,line width=1.35pt] (0.95,13)--(2.80,13);

    \draw[flow] (4.75,20.35)--(4.75,19.95);
    \draw[flow] (4.75,10.25)--(4.75,9.55);

    \outbox{mBase}
    \node[font=\scriptsize,text=vgInk,align=center] at (4.75,\Yeq)
      {$\hat y_t=\tilde y_t^{\mathrm{lin}}$};
    \miniout{\ob}
    % only interpolation shown in output
    \draw[mBase,line width=1.6pt](2.9,\ob+1.2)--(6.5,\ob+2.5);
    \foreach \s in {4,5,6}{
      \fill[mBase](0.2+\s*0.9,{\ob+1.2+(\s-3)*0.325})circle(0.10);
    }
  \end{scope}
\end{tikzpicture}
